% Формализация FrameNet

\begin{SCn}
\begin{small}

\scnheader{FrameNet{}}
\scnidtf{Язык формализации}
\scnidtf{FrameNet}
\scniselement{фреймовый язык}

\begin{scnrelfromlist}{разработчики}
    \scnitem{Чарльз Филлмор}
\end{scnrelfromlist}
\scnrelfrom{год создания}{1997}
\scnrelfrom{основная парадигма}{фреймовая семантика и корпусная лингвистика}

\scntext{описание}{Формальный язык представления знаний, основанный на фреймовой семантике. В его основе лежит идея, что значение слова или высказывания можно описать через фреймы — структурированные сценарии, отражающие типовые ситуации, события или отношения}

\scntext{назначение}{Служит для структурированного представления смысла в естественном языке. Его ключевая задача — формализовать связь между словами и стоящими за ними концепциями через унифицированные шаблоны — фреймы}

\begin{scnrelfromset}{основные конструкции}
    \scnfileitem{фреймы}
    \scnfileitem{семантические роли}
    \scnfileitem{лексические единицы}
    \scnfileitem{аннотированные предложения}
    \scnfileitem{семантические графы}
\end{scnrelfromset}

\scnrelfrom{стратегия вывода}{Контекстно-зависимый вывод на основе активации фреймов, иерархии фреймов, связей между фреймами}

\begin{scnrelfromlist}{поддержка}
    \scnfileitem{семантический анализ}
    \scnfileitem{разрешение неоднозначности}
    \scnfileitem{генерация текста}
\end{scnrelfromlist}

\begin{scnrelfromlist}{применение}
    \scnitem{вопросно-ответные системы}
    \scnitem{представление знаний}
    \scnitem{обработка естесственного языка}
    \scnitem{обучение ИИ-моделей}
    \scnitem{лингвистические исследования}
    \scnitem{машинный перевод}
\end{scnrelfromlist}

\begin{scnrelfromset}{примеры использования}
    \scnfileitem{Создание баз знаний, связывающих симптомы, заболевания и методы лечения, с возможностью логического вывода на основе здравого смысла}
    \scnfileitem{Реализация чат-ботов или виртуальных помощников, способных понимать контекст запросов и давать логически обоснованные ответы}
    \scnfileitem{Анализ тональности и выявление скрытых смыслов в документах}
    \scnfileitem{Аннотирование данных, что улучшает обучение моделей в задачах вроде распознавания образов или прогнозирования поведения пользователей}
    \scnfileitem{Создание интеллектуальных тренажёров, которые анализируют ошибки ученика и предлагают персонализированные пути обучения, опираясь на здравый смысл}
\end{scnrelfromset}

\begin{scnrelfromset}{реализации}
    \scnitem{Berkeley}
    \scnitem{Spanish FrameNet}
\end{scnrelfromset}

\begin{scnrelfromset}{особенности}
    \scnfileitem{гибкость}
    \scnfileitem{корпусная основа}
    \scnfileitem{интеграция с NLP}
\end{scnrelfromset}

\begin{scnrelfromset}{ограничения}
    \scnfileitem{языковая зависимость}
    \scnfileitem{сложность масштабирования}
    \scnfileitem{не поддерживает формальную логику}
\end{scnrelfromset}

\begin{scnrelfromset}{библиографические источники}
    \scnfileitem{Official FrameNet Website // International Computer Science Institute, Berkeley. – Электронный ресурс. Дата обращения: 27.05.2025. https://framenet.icsi.berkeley.edu/}
    \scnfileitem{FrameNet // Википедия. – 2025. – Электронный ресурс. Дата обращения: 27.05.2025. https://en.wikipedia.org/wiki/FrameNet}
    \scnfileitem{FrameNet: Sample usage // NLTK Documentation. – Электронный ресурс. Дата обращения: 27.05.2025. https://www.nltk.org/howto/framenet.html}
\end{scnrelfromset}

%--------------------------------------------------

\scnheader{KL-ONE}
\scnidtf{Язык представления знаний}
\scnidtf{KL-ONE}
\scniselement{фреймовый язык}

\begin{scnrelfromlist}{разработчики}
    \scnitem{Рональд Брахман}
    \scnitem{Джеймс Шмольце}
\end{scnrelfromlist}
\scnrelfrom{год создания}{1970-1980}
\scnrelfrom{основная парадигма}{структурное моделирование знаний с использованием иерархий концептов и ролей}

\scntext{описание}{Формальный язык представления знаний. Основная инновация — использование дедуктивного классификатора для автоматического вывода новых знаний на основе структурированной онтологии}

\scntext{назначение}{Преодоление семантической нечеткости в сетевых представлениях знаний через явное структурирование концептуальной информации в виде иерархии наследования. Служит для создания строго типизированных онтологий предметных областей}

\begin{scnrelfromset}{основные конструкции}
    \scnfileitem{концепты}
    \scnfileitem{роли}
    \scnfileitem{примитивные и определенные концепты}
    \scnfileitem{RoleSet}
\end{scnrelfromset}

\scnrelfrom{стратегия вывода}{Автоматизированный дедуктивный вывод на основе классификатора, проверяющего согласованность онтологии и выводящего новые отношения через анализ иерархии концептов и ролей}

\begin{scnrelfromlist}{поддержка}
    \scnfileitem{множественное наследование}
    \scnfileitem{проверка типов ролей}
    \scnfileitem{классификация концептов}
    \scnfileitem{выявление противоречий в онтологии}
\end{scnrelfromlist}

\begin{scnrelfromlist}{применение}
    \scnitem{экспертные системы}
    \scnitem{семантический анализ естественного языка}
    \scnitem{моделирование предметных областей}
    \scnitem{интеграция гетерогенных баз знаний}
\end{scnrelfromlist}

\begin{scnrelfromset}{примеры использования}
    \scnfileitem{Построение медицинских онтологий для диагностических систем, где заболевания классифицируются как концепты с симптомами в качестве ролей}
    \scnfileitem{Моделирование организационных структур предприятий с четким определением подчинённости отделов и должностей}
    \scnfileitem{Создание лингвистических ресурсов для обработки естественного языка, где слова связаны через концепты}
\end{scnrelfromset}

\begin{scnrelfromset}{реализации}
    \scnitem{KL-TWO}
    \scnitem{Krypton – комбинация KL-ONE с логическим языком}
    \scnitem{Loom}
    \scnitem{Ontolingua}
\end{scnrelfromset}

\begin{scnrelfromset}{особенности}
    \scnfileitem{строгая типизация концептов и ролей}
    \scnfileitem{формальная семантика}
    \scnfileitem{поддержка автоматизированного вывода}
    \scnfileitem{интеграция с логическими языками}
\end{scnrelfromset}

\begin{scnrelfromset}{ограничения}
    \scnfileitem{высокий порог входа из-за сложности формализма}
    \scnfileitem{ограниченная выразительность для некоторых типов знаний}
    \scnfileitem{трудности масштабирования для больших онтологий}
\end{scnrelfromset}

\begin{scnrelfromset}{библиографические источники}
    \scnfileitem{KL-ONE // Wikipedia. – 2024. – Электронный ресурс. Дата обращения: 27.05.2025. https://en.wikipedia.org/wiki/KL-ONE}
    \scnfileitem{Обзор языков представления знаний // Хабр. – 2020. – Электронный ресурс. Дата обращения: 27.05.2025. https://habr.com/ru/articles/526410/}
\end{scnrelfromset}

%=======================================

\scnheader{Сравнение фреймового языка KL-ONE и лингвистического ресурса FrameNet}
\begin{scneqtovector}
    \scnitem{KL-ONE}
    \scnitem{FrameNet}
\end{scneqtovector}

\begin{scnrelfromset}{сходства}
    \scnfileitem{Оба используют фреймовую модель представления знаний}
    \scnfileitem{Применяются для структурирования сложных концептуальных систем}
    \scnfileitem{Поддерживают иерархические отношения между элементами (наследование фреймов)}
    \scnfileitem{Используются в задачах искусственного интеллекта и обработки естественного языка}
\end{scnrelfromset}

\begin{scnrelfromset}{различия}
    \scnfileitem{KL-ONE — формальный язык логического вывода, FrameNet — лингвистическая база данных}
    \scnfileitem{KL-ONE использует строгую типизацию концептов и ролей, FrameNet допускает контекстную вариативность}
    \scnfileitem{KL-ONE ориентирован на автоматизированную классификацию, FrameNet — на аннотацию текстовых примеров}
    \scnfileitem{KL-ONE поддерживает дедуктивный вывод, FrameNet — семантический анализ}
    \scnfileitem{KL-ONE разработан для инженерии знаний, FrameNet — для лингвистических исследований и NLP}
\end{scnrelfromset}

\begin{scnrelfromset}{рекомендации по выбору}
    \scnfileitem{KL-ONE рекомендуется для задач, требующих строгой формализации онтологий и автоматического логического вывода}
    \scnfileitem{KL-ONE подходит для проектов с четкой структурой предметной области (медицина, телекоммуникации)}
    \scnfileitem{FrameNet целесообразно использовать для обработки естественного языка и семантического анализа текстов}
    \scnfileitem{FrameNet оптимален для лингвистических исследований, машинного перевода и вопросно-ответных систем}
    \scnfileitem{Для комплексных проектов рекомендуется комбинация: KL-ONE для онтологии, FrameNet для лингвистического интерфейса}
\end{scnrelfromset}

\end{small}
\end{SCn}